% =============================================================
%  0 引言
% =============================================================
\introduction

大数据和云计算的不断进步为物联网\textsuperscript{[1-2]}（Internet of
Things, IoT）提供了更强大的数据处理和存储能力，但云计算固有的时延并不能满足
实时性智能应用的服务质量需求\textsuperscript{[3]}。作为云计算的延伸和补充，
多接入边缘计算\textsuperscript{[4-5]}（Multi-Access Edge Computing, MEC）技术
通过就近为物联网设备提供算力和缓存资源，使得物联网设备能够执行计算密集型和
资源密集型应用，从而降低任务处理的时延以及能量的消耗\textsuperscript{[6-7]}。

在不同场景下部署 MEC 服务器，可以为物联网设备提供高效的任务处理服务
\textsuperscript{[8]}。将 MEC 与物联网融合\textsuperscript{[9]}，通过处理、
查询和存储多源物联网应用的数据流，有助于解决物联网应用中实时通信的需求问题。
鉴于物联网中对时延敏感和计算密集型应用需求的不断增长，文献\textsuperscript{[10]}
提出了缓存辅助的 MEC 卸载框架，通过联合优化通信、计算和缓存资源分配，
在动态缓存容量和计算资源受限的条件下最小化系统时延。文献\textsuperscript{[11]}
提出了基于李雅普诺夫优化理论的在线卸载和资源分配算法，解决了能效与服务时延之间的
平衡问题。上述研究充分发挥了 MEC 技术的优势，但未考虑在计算卸载过程中，
来自不同运营的实体之间缺乏信任的问题。

区块链作为一种去中心化且不可篡改的技术，提供了一个公开透明的公共账本，
能够有效解决边缘计算系统数据传输和处理时的信任缺乏问题\textsuperscript{[12]}。
将 MEC 集成到基于区块链的物联网系统中，以增强物联网设备的计算能力并保障所处理
任务的私密性\textsuperscript{[13]}。文献\textsuperscript{[14]}提出了一种区块链
支持的移动边缘计算平台，充分利用区块链技术的管理和控制功能，提高了网络中
物联网设备的服务质量（Quality of Service, QoS）。文献\textsuperscript{[15]}
研究了基于区块链的 MEC 系统中卸载决策、资源和区块间隔的联合优化问题，
以最大化区块链交易效率和处理速率。上述研究主要关注信任缺失问题。在基于区块链的
MEC 子系统中，边缘计算节点不仅为任务执行提供算力，还充当区块链节点，
区块的共识过程需要消耗大量的算力。因此，在任务执行和区块链共识之间进行合理有效的
计算资源分配，是一个关键的问题。

针对物联网设备在处理计算密集型及能耗密集型任务时能耗过高，以及来自不同运营机构的
边缘节点为物联网设备提供算力时缺乏信任等问题，拟提出一种区块链驱动的边缘计算系统
联合优化算法。首先，构建基于 MEC 和区块链的多基站网络架构，在各个基站部署边缘服务器，
为 IoT 设备就近提供算力。这些基站同时充当区块链节点，为边缘计算提供可信的环境。
通过将每个基站的算力在任务执行及区块共识之间进行合理分配，以及 IoT 设备的接入控制，
对所有 IoT 设备的平均总能耗之和最小化。其次，考虑到网络的动态特性和信息获取的不确定性，
将问题重建为马尔可夫决策过程（Markov Decision Process, MDP），并采用基于深度确定性
策略梯度（Deep Deterministic Policy Gradient, DDPG）的自适应决策算法进行求解。
DDPG 是一种专为解决具有连续变量问题而设计的深度强化学习（Deep Reinforcement
Learning, DRL）算法，允许在连续状态和动作空间中学习最优策略，而不依赖于显式的状态
转移概率。最后，将所提算法与随机接入控制算法和演员评论家（Actor-Critic, AC）算法
在能耗和系统累计奖励方面进行对比，验证所提算法的性能。
